%-------------------------------%
%  Author: Alessandro Sciarra   %
%    Date: 02 Jul 2025          %
%-------------------------------%

%===================================================================================================%
\subsection{Tests goals -- properties -- good principles}
%~~~~~~~~~~~~~~~~~~~~~~~~~~~~~~~~~~~~~~~~~~~~%
\begin{frame}{General idea and overview}
    \begin{overlayarea}{\textwidth}{0.7\textheight}
        \begin{center}
            \large
            \only<1-2>{There is slightly different advice around, pick your favourite:}
            \only<3->{There is one not-negotiable \textbf{requirement of tests}, though:}
        \end{center}
        \begin{varblock}{example}[0.8\textwidth]{Write F.I.R.S.T tests first}<only@1>
            \begin{description}[Self-validating:]
                \item[\PQ{F}ast:] To be run frequently
                \item[\PQ{I}ndependent:] To make failures meaningful
                \item[\PQ{R}epeatable:] To be run everywhere with same outcome
                \item[\PQ{S}elf-validating:] To be automatised
                \item[\PQ{T}imely:] To make production code testable
            \end{description}
        \end{varblock}
        \begin{varblock}{example}[\textwidth]{Google's perspective}<only@2>
            \begin{description}[Self-validating:]
                \item[Readable:] Tests are correct by inspection
                \item[Correct:] Do not rely on known bugs and test real scenarios
                \item[Complete:] Test most edge cases
                \item[Documenting:] Demonstration of how the API works
                \item[Resilient:] Repeatable, independent, hard to break, hermetic, etc.
            \end{description}
        \end{varblock}
        \begin{tikzpicture}[overlay, remember picture]
            \node[starburst, starburst points=21, starburst point height=15mm,
            minimum width=10cm, minimum height=6cm, Shape={red}{yellow!50}{red},
            line width=1.5pt, font=\Huge\bfseries, visible on=<4>]
            at ($(current page.center)-(0,1)$) {Existence};
        \end{tikzpicture}
    \end{overlayarea}
\end{frame}
%~~~~~~~~~~~~~~~~~~~~~~~~~~~~~~~~~~~~~~~~~~~~%
\begin{frame}[fragile]{Having tests is not enough, you must trust them!}{Few words about tests correctness}
    \begin{overlayarea}{\textwidth}{0.7\textheight}
        \begin{itemize}
            \item Never rely on known bugs
            \item<only@4-> Never mock the code that you are testing
        \end{itemize}
        \begin{varblock}{alert}[0.6\textwidth]{Having tests passing is simple...}<only@1-2>
            \begin{uncoverenv}<2->
                \begin{lstlisting}[style=MyCpp]
                    int cube(int x) {
                        //TODO: To be implemented
                        return 0;
                    }

                    TEST(cube_test) {
                        VERIFY(0 == cube(2));
                        VERIFY(0 == cube(3));
                        VERIFY(0 == cube(42));
                    }
                \end{lstlisting}
            \end{uncoverenv}
        \end{varblock}
        \begin{varblock}{example}[0.6\textwidth]{...better to have them failing!}<only@3>
            \begin{lstlisting}[style=MyCpp]
                int cube(int x) {
                    //TODO: To be implemented
                    return 0;
                }

                TEST(cube_test) {
                    VERIFY(8 == cube(2));
                    VERIFY(27 == cube(3));
                    VERIFY(74088 == cube(42));
                }
            \end{lstlisting}
        \end{varblock}
        \begin{varblock}{alert}[0.8\textwidth]{This is probably useless}<only@4>
            \begin{lstlisting}[style=MyCpp]
                class MockCoffeeMachine : public CoffeeMachine {
                    //For simplicity we assume always in temperature
                    double warmUp() override {
                        return constants::brewingT;
                    }
                }

                BOOST_AUTO_TEST_CASE(warmUp_test) {
                    MockCoffeeMachine machine{};
                    BOOST_REQUIRE_CLOSE(machine.warmUp(),
                    constants::brewingT, 0.1);
                }
            \end{lstlisting}
        \end{varblock}
        \vspace{3mm}
        \begin{varblock}{}[0.85\textwidth]{Think of a clear reaction procedure in your team}<only@5->
            What should happen in your project when a test fails?\\
            ...and when a bug that was not caught by tests is found?\\
            $\downarrow$\\
            \PB{Never miss any chance to improve your tests!}
        \end{varblock}
    \end{overlayarea}
\end{frame}
%~~~~~~~~~~~~~~~~~~~~~~~~~~~~~~~~~~~~~~~~~~~~%
\begin{frame}{Tests readability}
    \begin{center}
        \begin{tikzpicture}[every node/.style={inner sep=0mm, text=BGLIGHT}]
            \node (fig) {
\includegraphics[width=0.75\textwidth]{Balance}};
            \begin{DrawOnNode}{fig}
                \node[anchor=east] at (0.41,0.78) {
                    \begin{tabular}{c}
                        Avoid boilerplate\\
                        \footnotesize E.g. use functions for setup\\
                    \end{tabular}
                };
                \node[anchor=west] at (0.51,0.78) {
                    \begin{tabular}{c}
                        Give enough context\\
                        \footnotesize E.g. Not a single function for setup\\
                    \end{tabular}
                };
            \end{DrawOnNode}
        \end{tikzpicture}
    \end{center}
    \begin{varblock}{}[0.85\textwidth]{}
        \Large Remember: \PB{Tests are correct by inspection!}
    \end{varblock}
\end{frame}
%~~~~~~~~~~~~~~~~~~~~~~~~~~~~~~~~~~~~~~~~~~~~%
\begin{frame}[fragile]{Tests completeness: What should I test and how?}
    \begin{onlyenv}<1>
        \begin{itemize}
            \item Values  \\ \centerline{\includegraphics[scale=0.1]{Testing-value}}
            \item Actions \\ \centerline{\includegraphics[scale=0.12, clip, trim=0 4cm 0 4cm]{Testing-action}}
        \end{itemize}
        \vspace{3mm}
        \hfill\ldots{}and anything that makes sense!
    \end{onlyenv}
    \begin{varblock}{example}[0.98\textwidth]{A very nice example from Titus Winter}<only@2->
        \begin{overlayarea}{\textwidth}{0.55\textheight}
            \begin{uncoverenv}<3->
                \begin{lstlisting}[style=MyCpp]
                    int Factorial(int n) {
                        if (n == 1) return 1;
                        if (n == 5) return 120;
                        return -1; // TODO(goofus): figure this out.
                    }
                \end{lstlisting}
            \end{uncoverenv}
            \begin{onlyenv}<2-3>
                \begin{lstlisting}[style=MyCpp]
                    TEST(FactorialTest, BasicTests) {
                        EXPECT_EQ(1, Factorial(1));
                        EXPECT_EQ(120, Factorial(5));
                    }
                \end{lstlisting}
            \end{onlyenv}
            \begin{onlyenv}<4>
                \begin{lstlisting}[style=MyCpp]
                    TEST(FactorialTest, BasicTests) {
                        EXPECT_EQ(1, Factorial(1));
                        EXPECT_EQ(120, Factorial(5));
                        EXPECT_EQ(1, Factorial(0));
                        EXPECT_EQ(479001600, Factorial(12));
                        EXPECT_EQ(std::numeric_limits<int>::max() , Factorial(13));
                        EXPECT_EQ(1, Factorial(0));
                        EXPECT_EQ(std::numeric_limits<int>::max() , Factorial(-10));
                    }
                \end{lstlisting}
            \end{onlyenv}
        \end{overlayarea}
    \end{varblock}
    \begin{uncoverenv}<4>
        \begin{center}
            This will pretty much naturally emerge when doing TDD
        \end{center}
    \end{uncoverenv}
\end{frame}
%~~~~~~~~~~~~~~~~~~~~~~~~~~~~~~~~~~~~~~~~~~~~%
\begin{frame}{The hidden value of tests}{Few words about tests as API usage examples}
    \begin{overlayarea}{\textwidth}{0.68\textheight}
        \begin{center}
            \small
            Of course, a meme, do not take it (too) seriously, but...\\
            \vspace{3mm}
            \begin{uncoverenv}<2>
                
\includegraphics[width=0.7\textwidth]{Documentation-meme}\\
                \vspace{3mm}
                ...good tests should show how to use your code!
            \end{uncoverenv}
        \end{center}
        \FrameRemark<2>{In this sense, it should be avoided to use implementation details in tests (cf. white box testing downsides).}
    \end{overlayarea}
\end{frame}
%~~~~~~~~~~~~~~~~~~~~~~~~~~~~~~~~~~~~~~~~~~~~%
\begin{frame}[fragile]{Tests resilience: Google advice}
    \begin{onlyenv}<1>
        \begin{itemize}
            \setlength{\itemsep}{2mm}
            \item No flaky tests $\to$ Tests should always give the same outcome
            \item No brittle tests $\to$ One failure should not trigger many failures
            \item Tests reference values should not come from the system under test!
            \item Changing tests order should never change the outcome
            \item Tests should be as hermetic as possible $\to$ No I/O, no network, etc.
            \item No deep dependence $\to$ Avoid any implicit assumption
        \end{itemize}
        \vspace{3mm}
        \begin{varblock}{example}[0.5\textwidth]{}
            \large\bfseries\PS{Let's see some examples!}
        \end{varblock}
    \end{onlyenv}
    \begin{onlyenv}<2-3>
        \begin{varblock}{alert}[\textwidth]{A sneaky flaky test: What happens if the test-system is overloaded?}
            \begin{lstlisting}[style=MyCpp]
                TEST(UpdaterTest, RunsFast) {
                    Updater updater;
                    updater.UpdateAsync();
                    SleepFor(Seconds(.5)); // Half a second should be *plenty*.
                    EXPECT_TRUE(updater.Updated());
                }
            \end{lstlisting}
        \end{varblock}
        \begin{varblock}{alert}[\textwidth]{An infamous brittle pattern}<uncover@3>
            \begin{lstlisting}[style=MyCpp]
                TEST(Logger, LogWasCalled) {
                    StartLogCapture();
                    EXPECT_TRUE(Frobber::Start());
                    EXPECT_THAT(
                      Logs(), Contains("file.cc:421: Opened file frobber.config" )
                    );@|\tikzmark{logger}|@
                }
            \end{lstlisting}
        \end{varblock}
        \begin{tikzpicture}[remember picture, overlay]
            \draw[red, thick, double, visible on=<3>]
                ($(logger)+(3.05,0.12)$) -- ++(1.15,0)
                +(0.18,0) -- ++(0.65,0)
                +(2.4,0) -- ++(4.8,0);
        \end{tikzpicture}
    \end{onlyenv}
    \begin{onlyenv}<4-5>
        \begin{varblock}{alert}[\textwidth]{Where are the reference values coming from?}
            \begin{lstlisting}[style=MyCpp]
                BOOST_AUTO_TEST_CASE(Dirac_M)
                {
                    const hmc_float refs[4] =
                        { 2610.3804893063798, 4356.332327032359,
                          2614.2685771909237, 4364.1408252701831 };
                    test_fermionmatrix<physics::fermionmatrix::Dslash>(refs);
                }
            \end{lstlisting}
        \end{varblock}
        \begin{varblock}{alert}[\textwidth]{A non-hermetic test: What if run twice in parallel?}<uncover@5>
            \begin{lstlisting}[style=MyCpp]
                TEST(Server, StorageTest) {
                    StorageServer* server = GetStorageServerHandle();
                    auto my_val = rand();
                    server->Store("testkey", my_val);
                    EXPECT_EQ(my_val, server->Load("testkey"));
                }
            \end{lstlisting}
        \end{varblock}
    \end{onlyenv}
    \begin{onlyenv}<6>
        \vspace{-3mm}
        \begin{varblock}{alert}[\textwidth]{Deep dependence example}<only@6>
            \begin{lstlisting}[style=MyCpp]
                class File {
                    public:
                    // ...
                    virtual bool StatWithOptions(Stat_t* stat, StatOptions opts) {
                        return this->Stat(stat); // In base class ignore options
                    }
                };
                TEST(Filesystem, FSUsage) {
                    // ... and call to StatWithOptions
                    EXPECT_CALL(file, Stat(_)).Times(1);
                }  // This relies on a call to base StatWithOptions!!
            \end{lstlisting}
        \end{varblock}
        \begin{varblock}{quote}[\textwidth]{The law of implicit interfaces (AKA Hyrum's law)}[Hyrum Wright]<only@6>
            Given enough users of your public interface, soon or later someone will start implicitly relying on your implementation details (speed, memory consumption, etc.).\\[-1.2ex]~
        \end{varblock}
    \end{onlyenv}
    \begin{uncoverenv}<1-3,5->
        \FrameRemark{This list and all GoogleTest examples are taken from Titus Winters' and Hyrum Wright's talk at CppCon2015.}
    \end{uncoverenv}
\end{frame}
%~~~~~~~~~~~~~~~~~~~~~~~~~~~~~~~~~~~~~~~~~~~~%
\begin{frame}{A slide about tests precision}
    \vspace{-3mm}
    \begin{onlyenv}<1>
        \begin{center}
            
\includegraphics[width=0.7\textwidth]{Precision-meme}
        \end{center}
    \end{onlyenv}
    \begin{onlyenv}<2>
        \begin{itemize}
            \item We all know that machine precision is finite!
            \item Comparing floating point numbers has to be done carefully
            \item If you know a reference value exactly, put in the code more than the needed digits
            \PB{$\;\to\pi=3.141592653589793238462643$}
            \item Be aware of possible compiler techniques, e.g.\ FMA \Remark{fused multiply-add}\\
            \then This could make your test fail if you require too high precision
            \item Requiring too high precision makes tests brittle
            \item Requiring too low precision makes tests useless
            \item \alert{After all it is a judgement call!}
        \end{itemize}
    \end{onlyenv}
\end{frame}
%~~~~~~~~~~~~~~~~~~~~~~~~~~~~~~~~~~~~~~~~~~~~%
\begin{frame}{Few more words about reproducibility}
    \vspace{-5mm}
    \begin{varblock}{example}[0.96\textwidth]{How can I make my tests reproducible?}
        \begin{enumerate}
            \item Isolate them from non-deterministic aspects \Remark{e.g.\ network, database, random numbers}\\
            \then This is not always possible!
            \item Measure unwanted environment effect and subtract it in tests
            \item Detect if test can be run in the present environment\\
            \then If not, either fail or handle situation
            \item \alert{Handle intrinsic noise by statistics}\\
            \Then{E.g.\ if random numbers are needed, measure the failure frequency and re-run the test a number of times large enough to consider a failure as such!}
        \end{enumerate}
    \end{varblock}
\end{frame}
%~~~~~~~~~~~~~~~~~~~~~~~~~~~~~~~~~~~~~~~~~~~~%
\begin{frame}{Tests accuracy and reacting to failures}
    \begin{center}
        \begin{tikzpicture}
            \coordinate (O) at (0,0);
            \begin{scope}[every node/.style={inner ysep=6mm, inner xsep=3mm, font=\huge, minimum width=46mm},
                every label/.style={label distance=-3mm, font=\small, anchor=south, inner sep=0mm}]
                \node[anchor=south east, fill=PP!30, visible on=<2->] (M11) at (O) {
                    \begin{tabular}{c}
                        \emoji{top-arrow}\\[-5mm]
                        {\normalsize Ship it!}\\[-9pt]
                    \end{tabular}
                };
                \node[anchor=south west, fill=ALERT, visible on=<3->] (M12) at (O) {
                    \begin{tabular}{c}
                        \emoji{warning}\\[-5mm]
                        {\normalsize Add tests, fix code!}\\[-9pt]
                    \end{tabular}
                };
                \node[anchor=north east, fill=PT!10, visible on=<4->] (M21) at (O) {
                    \begin{tabular}{c}
                        \emoji{sweat-smile}\\[-5mm]
                        {\normalsize False alarm, fix tests!}\\[-9pt]
                    \end{tabular}
                };
                \node[anchor=north west, fill=PS!30, visible on=<5->] (M22) at (O) {
                    \begin{tabular}{c}
                        \emoji{tada}\\[-5mm]
                        {\normalsize Success story!}\\[-9pt]
                    \end{tabular}
                };
            \end{scope}
            \begin{scope}
                \draw (M11.south west) -- (M11.north west) -- (M12.north east) -- (M22.south east) --
                (M21.south west) -- (M21.north west) -- (M12.south east)
                (M11.north east) -- (M21.south east);
            \end{scope}
            \begin{scope}[node distance=2mm, every node/.style={inner sep=1mm, font=\Large}]
                \node[above = of M11] (CCyes) {\yes};
                \node[above = of M12] (CCno)  {\no};
                \node[left  = of M11, yshift=+3mm] (TPyes) {\yes};
                \node[left  = of M21, yshift=-3mm] (TPno)  {\no};
            \end{scope}
            \begin{scope}[every node/.style={font=\small}]
                \node at ($(CCyes.east)!0.5!(CCno.west)+(0,0)$) {Is code correct?};
                \node[rotate=90] at ($(TPyes.south)!0.5!(TPno.north)-(0,0)$) {Did tests pass?};
            \end{scope}
        \end{tikzpicture}
    \end{center}
    \FrameRemark{Tests accuracy means how much tests outcome matches code correctness.}
\end{frame}
%===================================================================================================%